% TOP_PROBLEM: {"id": 1487, "title": "Zeeman Conjecture", "category_id": 7, "category": {"id": 7, "name": "topology", "display_name": "Topology", "description": "Properties preserved under continuous deformations.", "slug": "topology", "order_index": 7, "created_at": "2026-07-31T15:26:25.671Z"}, "status": "open"}
% FIRST_POSTED: 2026-10-01
% ATTEMPT_STATUS: unresolved
% Imported from: unsolvedmath_additional_500.tex; UnsolvedMath ID 1487
% !TeX program = lualatex
% Compile twice: lualatex 1487.tex
% Self-contained: no external bibliography, images, or \input files.
% Numeric section labels and visible numbers are original UnsolvedMath IDs.
\documentclass[11pt,a4paper]{article}
\usepackage[margin=24mm,headheight=14pt]{geometry}
\usepackage{fontspec}
\setmainfont{Latin Modern Roman}
\setsansfont{Latin Modern Sans}
\setmonofont{Latin Modern Mono}
\usepackage{amsmath,amssymb,mathtools}
\usepackage{unicode-math}
\setmathfont{Latin Modern Math}
\usepackage{microtype}
\usepackage{enumitem,xurl,needspace}
\usepackage[dvipsnames]{xcolor}
\usepackage{hyperref}
\hypersetup{unicode=true,colorlinks=true,linkcolor=MidnightBlue,urlcolor=MidnightBlue,
 pdftitle={UnsolvedMath Problem 1487},pdfauthor={Source-annotated compilation},bookmarksnumbered=true}
\usepackage{bookmark,tocloft,titlesec,fancyhdr}
\setlength{\cftsecnumwidth}{4.2em}
\cftsetpnumwidth{2.5em}
\cftsetrmarg{4em}
\titleformat{\section}{\Large\bfseries\raggedright}{\thesection}{0.7em}{}
\pagestyle{fancy}\fancyhf{}
\fancyhead[L]{\small\sffamily Additional UnsolvedMath statements}
\fancyhead[R]{\small Original catalogue IDs}
\fancyfoot[C]{\thepage}
\setlength{\parindent}{0pt}
\setlength{\parskip}{5pt plus 1pt}
\setlength{\emergencystretch}{3em}
\setcounter{tocdepth}{1}\setcounter{secnumdepth}{1}
\setlist[itemize]{leftmargin=3em,itemsep=4pt,topsep=4pt,parsep=0pt}
\allowdisplaybreaks
\newcommand{\N}{\mathbb{N}}
\newcommand{\Z}{\mathbb{Z}}
\newcommand{\Q}{\mathbb{Q}}
\newcommand{\R}{\mathbb{R}}
\newcommand{\C}{\mathbb{C}}
\newcommand{\F}{\mathbb{F}}
\newcommand{\A}{\mathbb{A}}
\newcommand{\PP}{\mathbb{P}}
\newcommand{\E}{\mathbb{E}}
\newcommand{\eps}{\varepsilon}
\newcommand{\dd}{\,\mathrm d}
\newcommand{\sref}[2]{\hyperlink{source:#1:#2}{[S#2]}}
\newif\ifshowcatalogue
\showcataloguetrue
\newcommand{\cataloguescope}[1]{\ifshowcatalogue
 \paragraph{Statement recorded in the supplied source.}
 #1\fi}
\begin{document}
% UnsolvedMath ID 1487; source catalogue code TOP-009

\setcounter{section}{1486}

\section{Zeeman’s Interval-Product Collapse Conjecture}\label{1487}

{\small\sffamily UnsolvedMath ID 1487 · Topology}

\subsection{Definitions and mathematical statement}

\paragraph{Shared notation.}Unless a section says otherwise, $\N=\{1,2,\ldots\}$,
$\Z$ is the ring of integers, $\Q,\R,\C$ are the rational, real, and complex
fields, $[N]=\{1,\ldots,\lfloor N\rfloor\}$, and $\log$ is the natural
logarithm. For real functions of a parameter tending to infinity,
$f\ll g$ and $f=O(g)$ mean $|f|\leq Cg$ eventually for some fixed $C$;
$f\gg g$ reverses this comparison; $f=o(g)$ means $f/g\to0$;
$f\sim g$ means $f/g\to1$; and $f\asymp g$ means both order bounds.
A subscript on $O$ or $\ll$ specifies allowed constant dependence.
The expression $n^{o(1)}$ in an upper bound means growth smaller than
$n^\epsilon$ for each fixed $\epsilon>0$. Local definitions govern any
exception, finite-field convention, or use of the same symbol for another
object. An asymptotic claim is not automatically an effective algorithm.



A finite complex is contractible if its identity map is homotopic to a constant map. An elementary collapse removes a cell together with a free face, and collapsibility means a finite sequence of such removals reaches one vertex. The standard polyhedral formulation allows a suitable triangulation or subdivision of the product with an interval. Contractibility alone does not supply such a collapse sequence. \sref{1487}{1} \sref{1487}{2}

\paragraph{Statement recorded in the supplied source.}
Is $K \times [0,1]$ collapsible for every finite contractible 2-dimensional CW complex K? \sref{1487}{1}

\paragraph{Scope and interpretation.}

The supplied record says CW complex. The elementary-collapse convention is the standard finite polyhedral/simplicial one in the cited formulation; no claim is made for arbitrary nonregular cell attachments without specifying a compatible collapse model. \sref{1487}{1}

\subsection{Short English statement}

Does multiplying a finite contractible two-dimensional complex by an interval always make it collapsible?

\subsection{Research attempt: explicit product collapses and the contractibility gap}
\textbf{Status and category.} Work with finite simplicial complexes and their regular polyhedral product cells, allowing compatible subdivisions as in the selected formulation. We prove the conjectured conclusion when $K$ is already collapsible. We do not replace contractibility by collapsibility in the general question. The primary discussion [R1] records Zeeman's polyhedral conjecture and provides context for this distinction.

\subsubsection{1. Every finite cylinder collapses onto its bottom}
Give $K\times I$ the product cell structure, with cells $\sigma\times\{0\}$, $\sigma\times\{1\}$ and $\sigma\times(0,1)$ for simplices $\sigma$ of $K$. Process the simplices of $K$ in decreasing dimension. For each $\sigma$, remove the prism $\sigma\times I$ together with its top face $\sigma\times\{1\}$. At that stage every prism and top face belonging to a proper coface of $\sigma$ has already been removed. Hence the top face is contained in no remaining larger cell except its own prism, and is free. The bottom face is retained throughout. This gives
\[K\times I\searrow K\times\{0\}.\tag{1}\]
Product cells are balls attached along their boundaries; the standard relative triangulation of a free polyhedral pair realizes each of these collapses after compatible subdivision. This is the polyhedral collapse convention in the target, not an assertion about arbitrary nonregular CW cells.

\subsubsection{2. Collapses of the base can also be transported}
Suppose $K\searrow K'$ removes a simplex $\sigma$ and its free codimension-one face $\tau$. In the product, first remove $\sigma\times I$ paired with $\tau\times I$. Freeness follows because $\tau$ has no other coface in $K$. Then remove $\sigma\times\{0\}$ with $\tau\times\{0\}$, and remove the corresponding top pair. The result is exactly $K'\times I$. This provides a second constructive way of lifting a collapse sequence of the base.

In particular, if $K\searrow\{v\}$ then either this lifted sequence followed by collapsing the final interval, or (1) followed by the original sequence, proves $K\times I$ collapsible. Every finite tree is an elementary example: remove a leaf edge with its leaf vertex until one vertex remains. Products of these examples can be treated by iterating the same construction.

\subsubsection{3. Why the universal cylinder collapse does not settle the conjecture}
Equation (1) holds even for noncontractible $K$. For example, a circle times an interval collapses onto a circle, but it cannot collapse to a point because elementary collapses preserve homotopy type and its first homology is nonzero. Thus a collapse onto a subcomplex is not a collapse of that subcomplex.

For contractible $K$ the homology obstruction disappears, but (1) still ends at $K$ itself. If that chosen triangulation has no available free face, the exhibited sequence stops. The target asks whether another sequence, using the thickness of the cylinder in a different way and possibly a subdivision, can continue to a point. Choosing a particularly unhelpful collapse sequence is not a disproof either: collapsibility asks for existence of some sequence.

\subsubsection{4. Verification and exact remaining task}
The free-face audit in (1) depends on decreasing dimension; increasing dimension would leave coface prisms and make the proposed top faces nonfree. The transported collapse also removes the lateral prism pair before the top and bottom pairs, for the same reason. These local checks give explicit valid sequences in the stated subcase. No algorithm here converts a contraction homotopy of an arbitrary contractible 2-complex into such a sequence. The missing conversion, while preserving the allowed product and subdivision category, is exactly the general Zeeman question.

\subsection{Sources}

{\small

\begin{itemize}

\item[\hypertarget{source:1487:1}{S1}] ULAM.AI, \emph{UnsolvedMath}, supplied \texttt{problems.json}, record 1487 (TOP-009), statement and background. The embedded literature-review date is 2026-08-17; that is the archive’s review date, not a fresh certification. Dataset landing page: \url{https://huggingface.co/datasets/ulamai/UnsolvedMath}.

\item[\hypertarget{source:1487:2}{S2}] Zeeman's conjecture overview. \url{https://en.wikipedia.org/wiki/Zeeman%27s_conjecture}. Bibliographic information imported from the supplied record.

\item[R1] Zeeman’s Collapsing Conjecture, chapter XI of Two-Dimensional Homotopy and Combinatorial Group Theory, primary monograph discussion.
\url{https://www.cambridge.org/core/books/abs/twodimensional-homotopy-and-combinatorial-group-theory/zeemans-collapsing-conjecture/AE8D58F55223AF72C7DBB5CD88A6E9D6}.
\end{itemize}

}

\label{end:1487}
\end{document}
